%! AP Statistics — Unit 8, Lesson 8.1 — Group Activity (Answer Key)
\documentclass[10pt]{article}
\usepackage{apstats-article}
\usepackage{apstats-key}

\begin{document}

\pageheader{Unit 8, Lesson 8.1}{Group Activity --- Answer Key}
\namepartnerperiod

\vspace{0.1in}
\textbf{Are These Results Surprising?}

\vspace{0.05in}
A school surveys 120 students about their preferred study environment. The results are:

\vspace{0.08in}
\renewcommand{\arraystretch}{1.4}
\begin{tabularx}{0.7\linewidth}{Xcc}
    \textbf{Environment} & \textbf{Observed Count} & \textbf{Expected Count} \\
    \hline
    Library    & 38 & \ans{30} \\
    Home       & 52 & \ans{30} \\
    Caf\'e     & 18 & \ans{30} \\
    Classroom  & 12 & \ans{30} \\
    \hline
    \textbf{Total} & \textbf{120} & \ans{120} \\
\end{tabularx}

\vspace{0.1in}
\textbf{Directions:} Work with your partner to complete parts (a)--(d) below.

\begin{tcolorbox}[colback=white, colframe=black!40, title=\textbf{Tier R --- Remediate},
    fonttitle=\bfseries, arc=2mm, left=3mm, right=3mm, top=2mm, bottom=2mm]

\textbf{Part (a)} --- Fill in the expected count for each category, assuming all four
study environments are equally preferred.

\vspace{0.05in}
\textit{Hint:} If 120 students are equally spread among 4 categories, expected count
$= 120 \div 4 = $ \ans{30} for each category.

\vspace{0.1in}
Go back and fill in the Expected Count column in the table above.
\end{tcolorbox}

\begin{tcolorbox}[colback=white, colframe=black!40, title=\textbf{Tier A --- Approaching Proficiency},
    fonttitle=\bfseries, arc=2mm, left=3mm, right=3mm, top=2mm, bottom=2mm]

\textbf{Part (b)} --- For each category, compute Observed $-$ Expected. Record your
answers below.

\vspace{0.05in}
\begin{tabularx}{\linewidth}{Xcccc}
    & Library & Home & Caf\'e & Classroom \\
    \hline
    Observed $-$ Expected & \ans{$+8$} & \ans{$+22$} & \ans{$-12$} & \ans{$-18$} \\
\end{tabularx}

\vspace{0.1in}
Which category has the largest \textbf{positive} difference? \ans{Home ($+22$)}

Which category has the largest \textbf{negative} difference? \ans{Classroom ($-18$)}

\vspace{0.1in}
\textbf{Part (c)} --- Without doing any formal calculation, do you think the differences
are large enough to conclude that students do \emph{not} have equal preferences? Explain
your reasoning in 2--3 sentences.

\ansline{Answers vary --- accept reasonable judgment. The Home and Classroom differences are the largest (22 and 18 away from an expected 30), which may suggest real differences in preference. However, without a formal test we cannot rule out that these differences are due to chance alone.}

\begin{teachernote}
    Students should recognize that large differences (especially Home $+22$ and
    Classroom $-18$) are potentially meaningful but cannot be confirmed without a
    $p$-value. Accept a range of responses as long as students acknowledge uncertainty.
\end{teachernote}
\end{tcolorbox}

\begin{tcolorbox}[colback=white, colframe=black!40, title=\textbf{Tier E --- Extension},
    fonttitle=\bfseries, arc=2mm, left=3mm, right=3mm, top=2mm, bottom=2mm]

\textbf{Part (d)} --- What additional information would help you decide whether these
differences are \emph{statistically significant}? In other words, what would a formal
statistical test provide that your informal comparison in part (c) cannot?

\ansline{We would need to know the probability of seeing differences this large or larger by chance alone, assuming students have equal preferences. That is exactly what the chi-square test provides --- a $p$-value that quantifies how unusual the observed gaps are under the null hypothesis of equal preference.}

\textbf{Bonus:} The chi-square statistic adds up all the gaps in a specific way.
Why might we need to \emph{adjust} each gap relative to the expected count, rather than
just adding up the raw gaps?

\ansline{A gap of 10 from an expected count of 100 is less surprising than a gap of 10 from an expected count of 11. Dividing by the expected count scales each gap to reflect how unusual it is --- a large gap relative to a small expected count is more meaningful than the same raw gap relative to a large expected count.}
\end{tcolorbox}

\end{document}
